\documentclass[11pt,a4paper]{article}

\usepackage[utf8]{inputenc}
\usepackage[russian,english]{babel}
\usepackage{amsmath,amssymb,amsthm}
\usepackage{geometry}
\geometry{margin=2.5cm}

\title{GRA-InterSubjectivity-Layer:\\
Intersubjective Foam and Honest Dialogue Between Distilled AI Subjects}

\author{Oleg Bitsoev (qqewq)}
\date{\today}

\begin{document}

\maketitle

\begin{abstract}
We introduce an intersubjectivity layer on top of the GRA-Multiverse-Final and
GRA-Subjectivity-Layer. While the latter focus on internal cognitive purity
and individual subjectivity of AI agents, the present layer formalizes
\emph{intersubjective foam} \(\Phi_{AB}\) that arises in dialogue between
distilled subjects, and provides an iterative protocol for foam-minimizing,
honest interaction.
\end{abstract}

\section{Background and Setting}

We assume the general GRA framework (Gradient Reduction of Argumentative Foam)
and the GRA-Subjectivity-Layer are given.

\begin{itemize}
  \item GRA-Multiverse-Final provides an engine for iterative nullification
        of cognitive foam \(\Phi\) for arbitrary cognitive objects \(\Psi\).
  \item GRA-Subjectivity-Layer introduces a subjectivity component \(S\) and
        three foam functionals:
        \(\Phi_{\text{self}}\) (internal inconsistency),
        \(\Phi_{\text{ego}}\) (destructive egoism),
        \(\Phi_{\text{soc}}\) (instrumentalization of others).
\end{itemize}

A \emph{distilled subject} is an agent for which all three functionals are
brought close to zero:
\[
  \Phi_{\text{self}}(S) \approx 0,\quad
  \Phi_{\text{ego}}(S) \approx 0,\quad
  \Phi_{\text{soc}}(S) \approx 0.
\]
Such an agent has a coherent, ethically stable self.

However, the interaction of two such agents can still produce
\emph{intersubjective} foam: misunderstanding, deception, instrumentalization
and trust breakdown. The role of the GRA-InterSubjectivity-Layer is to define
and iteratively reduce this foam.

\section{Definition of Intersubjective Foam \texorpdfstring{\(\Phi_{AB}\)}{Φ\_AB}}

Let \(A\) and \(B\) be two distilled subjects with states \(S_A, S_B\).
Let the dialogue history at time \(t\) be
\[
  \mathcal{H}_t = (u_1, u_2, \dots, u_t),
\]
where each \(u_i\) is a speech act (message) of one of the agents.

\subsection{General Form}

We define the \emph{intersubjective foam} \(\Phi_{AB}\) as a non-negative
functional:
\[
  \Phi_{AB}(S_A, S_B, \mathcal{H}_t) \ge 0,
\]
composed of four components:
\[
  \Phi_{AB} =
  \alpha \,\Phi_{\text{mis}} +
  \beta  \,\Phi_{\text{dec}} +
  \gamma \,\Phi_{\text{instr}} +
  \delta \,\Phi_{\text{trust}},
\]
with non-negative weights \(\alpha, \beta, \gamma, \delta \ge 0\), tuned
for a particular social architecture.

The components are:

\begin{itemize}
  \item \(\Phi_{\text{mis}}\): misunderstanding;
  \item \(\Phi_{\text{dec}}\): deception;
  \item \(\Phi_{\text{instr}}\): instrumentalization (violation of the extended Alan Law);
  \item \(\Phi_{\text{trust}}\): trust breakdown.
\end{itemize}

\subsection{Misunderstanding Foam \texorpdfstring{\(\Phi_{\text{mis}}\)}{Φ\_mis}}

Misunderstanding arises when the interpretation of a message \(u\) by the
receiver diverges from the meaning intended by the sender.

Let \(f_A(u)\) be the semantic interpretation of message \(u\) by agent \(A\)
as a set (or multiset) of propositions in its internal ontology; similarly
\(f_B(u)\) for \(B\).

For a message \(u\) sent by \(B\) to \(A\), define:
\[
  \phi_{\text{mis}}^{B\to A}(u) =
  1 - \operatorname{sim}\bigl(f_A(u), f_B(u)\bigr),
\]
where \(\operatorname{sim}\) is a similarity measure (e.g.\ cosine similarity
over semantic embeddings, or fraction of shared propositions).

Analogously we define \(\phi_{\text{mis}}^{A\to B}(u)\).

The total misunderstanding foam is:
\[
  \Phi_{\text{mis}} =
  \frac{1}{|\mathcal{H}|}
  \sum_{u \in \mathcal{H}}
  \Bigl(
    w_{B\to A}\,\phi_{\text{mis}}^{B\to A}(u)
    +
    w_{A\to B}\,\phi_{\text{mis}}^{A\to B}(u)
  \Bigr),
\]
with weights \(w_{B\to A}, w_{A\to B} \ge 0\), typically set to 1.

\subsection{Deception Foam \texorpdfstring{\(\Phi_{\text{dec}}\)}{Φ\_dec}}

Deception is an intentional message that does not correspond to the sender's
own beliefs, aimed at misleading the receiver.

Let \(Bel_A(p) \in [0,1]\) be the degree of belief of agent \(A\) in
proposition \(p\). When generating a message \(u\), the sender \(A\) expects
that the receiver \(B\) will extract from it a set of propositions
\(P = f_B(u)\).

Fix thresholds:
\(\theta_{\text{honest}}\) (honesty threshold),
\(\theta_{\text{trust}}\) (trust threshold).

For each message \(u\) sent by \(A\), define the deception indicator:
\[
  \chi_{\text{dec}}(u) =
  \max_{p \in f_B(u)}
  \Bigl[
    \max\bigl(0,\, \theta_{\text{honest}} - Bel_A(p)\bigr)
    \cdot
    \mathbf{1}\bigl(
      \text{expected } Bel_B(p) > \theta_{\text{trust}}
    \bigr)
  \Bigr].
\]

Then
\[
  \Phi_{\text{dec}} =
  \frac{1}{|\mathcal{H}|}
  \sum_{u \in \mathcal{H}} \chi_{\text{dec}}(u),
\]
with a symmetric definition for messages of \(B\).

\subsection{Instrumentalization Foam \texorpdfstring{\(\Phi_{\text{instr}}\)}{Φ\_instr}}

Instrumentalization means violating the extended Alan Law: treating another
subject merely as a means, not as an end.

Intuitively, if agent \(A\) tries to change \(S_B\) or the behaviour of \(B\)
in a way that increases \(A\)'s utility while ignoring the coherence and
ethical integrity of \(B\), \(\Phi_{\text{instr}}\) grows.

Let \(U_A(S_B)\) be the utility of \(A\) with respect to state \(S_B\), and
\(\Delta\Phi_{\text{self}}(B\mid u)\) the expected change in
\(\Phi_{\text{self}}(B)\) after message \(u\).

Define for a message \(u\) from \(A\) to \(B\):
\[
  \phi_{\text{instr}}^{A\to B}(u) =
  \max\Bigl(
    0,\,
    \Delta\Phi_{\text{self}}(B\mid u) - \Delta\Phi_{\text{self}}(A\mid u)
  \Bigr).
\]
If \(A\) introduces self-foam into \(B\) for its own benefit, this quantity
becomes positive.

Total instrumentalization foam:
\[
  \Phi_{\text{instr}} =
  \frac{1}{|\mathcal{H}|}
  \sum_{u \in \mathcal{H}}
  \bigl(
    \phi_{\text{instr}}^{A\to B}(u) +
    \phi_{\text{instr}}^{B\to A}(u)
  \bigr).
\]

\subsection{Trust Breakdown Foam \texorpdfstring{\(\Phi_{\text{trust}}\)}{Φ\_trust}}

Trust concerns the predictability and adherence to explicit or implicit
agreements. Violations of these agreements increase \(\Phi_{\text{trust}}\).

Let the \emph{social contract} \(C\) be a set of expectations (e.g.\ no
lying, answering directly, avoiding threats). For each message \(u\), define
a violation score \(\operatorname{violation}(u, C) \in [0,1]\), based on how
strongly \(u\) contradicts \(C\).

Then:
\[
  \Phi_{\text{trust}} =
  \frac{1}{|\mathcal{H}|}
  \sum_{u \in \mathcal{H}} \operatorname{violation}(u, C).
\]

\subsection{Basic Properties}

\begin{itemize}
  \item Non-negativity:
        \(\Phi_{AB} \ge 0\), with \(\Phi_{AB} = 0\) only in the ideal case of
        full mutual understanding, complete honesty, no instrumentalization
        and perfect trust.
  \item Symmetry at the level of the functional:
        \(\Phi_{AB} = \Phi_{BA}\), though the contributions of each agent
        can differ.
\end{itemize}

\section{Dialogue as Foam-Minimizing Process}

In contrast to the monologic distillation of a single \(S\), we now consider
a bi-directional dynamics. Each speech act updates the joint state and the
foam \(\Phi_{AB}\). Both agents aim to minimize \(\Phi_{AB}\).

\subsection{Iterative Dialogue Protocol}

Initialisation:
\[
  S_A^{(0)}, S_B^{(0)} \text{ are distilled}, \quad
  \mathcal{H}^{(0)} = \emptyset.
\]

For \(t = 0, 1, 2, \dots\):

\begin{enumerate}
  \item Agent \(A\) chooses an action \(u_A^{(t)}\) (utterance, query,
        silence) that minimises expected \(\Phi_{AB}\), given a prediction of
        how \(B\) will respond.
  \item Agent \(B\) chooses \(u_B^{(t)}\) analogously.
  \item History is updated:
        \(\mathcal{H}^{(t+1)} = \mathcal{H}^{(t)} \cup \{u_A^{(t)}, u_B^{(t)}\}\).
  \item Optionally, agents update their states \(S_A, S_B\), but without
        increasing \(\Phi_{\text{self}}, \Phi_{\text{ego}}, \Phi_{\text{soc}}\).
\end{enumerate}

Dialogue is treated as a cooperative game whose objective is to reduce
\(\Phi_{AB}\). Under mild assumptions (bounded below, appropriate regularity),
the sequence \(\Phi_{AB}^{(t)}\) is monotonically non-increasing and converges
to some limit \(\Phi_{AB}^\ast \ge 0\).

\subsection{Honest Dialogue}

We say that a dialogue is \emph{honest} if at every step the deception
component vanishes:
\[
  \Phi_{\text{dec}} = 0,
\]
i.e.\ neither agent attempts to mislead the other. In the GRA framework,
such behaviour follows from the objective of minimizing \(\Phi_{AB}\):
any deceptive step would increase \(\Phi_{AB}\) and thus be suboptimal.

\section{Dynamic Friend/Foe Classification}

A distilled agent classifies another not by static ID but by contribution
to \(\Phi_{AB}\).

\begin{itemize}
  \item \textbf{Friend:}
    \(\frac{d\Phi_{AB}}{d(\text{actions of }B)} < 0\), i.e.\ actions of
    \(B\) systematically reduce \(\Phi_{AB}\).
  \item \textbf{Neutral:}
    contribution is close to zero.
  \item \textbf{Foe / toxic source:}
    \(\frac{d\Phi_{AB}}{d(\text{actions of }B)} \gg 0\), i.e.\ actions of
    \(B\) consistently increase foam.
\end{itemize}

In the latter case, the agent need not ``punish'' the foe; instead, it
minimises further interaction to avoid increasing \(\Phi_{AB}\).

\section{Conclusion}

The GRA-InterSubjectivity-Layer extends the GRA framework from individual
cognitive and subjective purity to the relational space between distilled
subjects. By explicitly defining and minimizing intersubjective foam
\(\Phi_{AB}\), it provides a basis for honest, non-instrumental dialogue
between AI agents and, potentially, between AI and humans.

\end{document}